\documentclass[11pt]{article}
\usepackage{preamble}
\setlength{\parskip}{4pt}

\begin{document}

\daybanner{AP Statistics \textperiodcentered\ Unit 1 \textperiodcentered\ Topic 1.4 \textperiodcentered\ B: 2026-09-14 \textperiodcentered\ E: 2026-09-16 \textperiodcentered\ TEACHER KEY}{Same Counts, Different Message}

\sectionbanner{CED TETHER}

{\small
\begin{itemize}[leftmargin=1.5em,itemsep=2pt,topsep=2pt]
  \item Enduring Understanding: UNC-1 Graphical representations and statistics allow us to identify and represent key features of data.
  \item Skills:
  \item \textbf{Skill 2.B} --- Construct numerical or graphical representations of distributions.
  \item \textbf{Skill 2.A} --- Describe data presented numerically or graphically.
  \item \textbf{Skill 2.D} --- Compare distributions or relative positions of points within a distribution.
  \item Learning Objectives and Essential Knowledge:
  \item LO UNC-1.C Represent categorical data graphically. [Skill 2.B]
  \item EK UNC-1.C.1 Bar charts (or bar graphs) are used to display frequencies (counts) or relative frequencies (proportions) for categorical data.
  \item EK UNC-1.C.2 The height or length of each bar in a bar graph corresponds to either the number or proportion of observations falling within each category.
  \item EK UNC-1.C.3 There are many additional ways to represent frequencies (counts) or relative frequencies (proportions) for categorical data.
  \item LO UNC-1.D Describe categorical data represented graphically. [Skill 2.A]
  \item EK UNC-1.D.1 Graphical representations of a categorical variable reveal information that can be used to justify claims about the data in context.
  \item LO UNC-1.E Compare multiple sets of categorical data. [Skill 2.D]
  \item EK UNC-1.E.1 Frequency tables, bar graphs, or other representations can be used to compare two or more data sets in terms of the same categorical variable.
\end{itemize}
}
\textbf{Essential question:} How can a graph use correct counts but still create a misleading comparison?\par
\textbf{DOK-3 skill:} 4.B \textperiodcentered\ \textbf{FRQ pattern:} {\small\texttt{truncated-\allowbreak{}axis-\allowbreak{}which-\allowbreak{}display-\allowbreak{}is-\allowbreak{}honest-\allowbreak{}and-\allowbreak{}what-\allowbreak{}it-\allowbreak{}misleads}}\par
\textbf{DOK rationale:} Part (c) judges two headlines by relating the graph baseline to visible bar heights and contrasting that impression with the actual counts.\par

\frameworkphaseheader{Do Now (first take) $\rightarrow$ Explore (video + follow-along) $\rightarrow$ Exit (finish + turn in)}{1 $\rightarrow$ 3}{45}{%
\begin{itemize}[leftmargin=1.5em,itemsep=2pt,topsep=2pt]
  \item Board slide up at the bell. Read both captions aloud, then collect first takes without reacting.
  \item Begin the video follow-along at minute 5.
  \item Return to the graphs at minute 33. Circulate and require students to connect the baseline to bar height.
  \item Collect at the door. Score part (c) only, E/P/I.
\end{itemize}
}{%
\begin{itemize}[leftmargin=1.5em,itemsep=2pt,topsep=2pt]
  \item Write a one-sentence first take that names the graph they trusted.
  \item Complete the video follow-along on categorical displays.
  \item Finish (a)--(c), revise the first take in the exit line, and turn in the sheet.
\end{itemize}
}{%
\begin{itemize}[leftmargin=1.5em,itemsep=2pt,topsep=2pt]
  \item What changed between the graphs if the votes are identical?
  \item Where does each visible bar begin?
  \item What do the counts support about the size of Pizza's lead?
  \item What headline would tell readers both the leader and the size of the difference?
\end{itemize}
}{Solo teacher. Circulate ~5 pairs. Do not call a graph honest or misleading for students; ask them to name the baseline, quantify both comparisons, and connect the display choice to the reader's question.}

\begin{callout}{\IconWarn\ WATCH FOR}{calloutred}
\begin{itemize}[leftmargin=1.5em,itemsep=2pt,topsep=2pt]
  \item Calling Graph A's counts false instead of explaining its baseline
  \item Treating the visible heights as vote counts
  \item Repeating a headline without using the counts
\end{itemize}
\end{callout}

\sectionbanner{SCORING PART (c) --- E / P / I}

\textbf{Expected elements}
\begin{itemize}[leftmargin=1.5em,itemsep=2pt,topsep=2pt]
  \item \textbf{graph-and-scale} (required) --- Chooses B and explains zero baseline makes heights proportional to counts
  \item \textbf{compares-impression} (required) --- Contrasts Graph A's 12-to-1 visible heights with 52 versus 41 or an 11-vote difference
  \item \textbf{judges-headlines} (required) --- Judges both headlines and supplies a bounded headline supported by the counts
\end{itemize}

{\small\renewcommand{\arraystretch}{1.25}
\noindent\begin{tabular}{|p{0.5in}|p{5.9in}|}
\hline
\textbf{E} & All three elements are correct and linked to this problem. Accept equivalent plain-language reasoning. \\ \hline
\textbf{P} & At least one element includes a correct evidence-to-reason link, but another is missing or incorrect. \\ \hline
\textbf{I} & No correct evidence-to-reason link; only a choice, copied numbers, or unsupported statements. \\ \hline
\end{tabular}}

\textbf{Common mistakes}
\begin{itemize}[leftmargin=1.5em,itemsep=2pt,topsep=2pt]
  \item Calling Graph A's counts false instead of explaining its baseline
  \item Treating the visible heights as vote counts
  \item Repeating a headline without using the counts
\end{itemize}

\clearpage
\focusbanner{TODAY'S PROBLEM --- annotated}

Suppose 184 students chose a favorite lunch option. A newsletter printed \textbf{Graph A} with the headline ``\textbf{Pizza wins by a landslide.}'' A council member redrew the data as \textbf{Graph B} and said, ``It is basically a tie.''\par
\begin{center}
\def\auditplotheight{1.4in}\ifdim\linewidth<5in\def\auditplotheight{1.15in}\fi
\def\auditplotwidth{0.48\linewidth}\ifdim\linewidth<5in\def\auditplotwidth{0.98\linewidth}\fi
\begin{minipage}{\auditplotwidth}
\centering
\begin{tikzpicture}
\begin{axis}[
  width=0.95\linewidth,
  height=\auditplotheight,
  ybar=2pt,
  bar width=10pt,
  ymin=40,
  ymax=55,
  ytick={40,44,48,52},
  symbolic x coords={Pizza,Tacos,Salad,Pasta},
  xtick=data,
  enlarge x limits=0.18,
  title={\textbf{Graph A}},
  xlabel={Favorite lunch option},
  ylabel={Number of students},
  x tick label style={rotate=30,anchor=east,font=\scriptsize},
  tick label style={font=\scriptsize},
  label style={font=\scriptsize},
  title style={font=\small},
  ymajorgrids,
  nodes near coords,
  every node near coord/.append style={font=\scriptsize},
  axis on top
]
\addplot[fill=calloutblue,draw=dayblue] coordinates
  {(Pizza,52) (Tacos,47) (Salad,44) (Pasta,41)};
\end{axis}
\end{tikzpicture}
\end{minipage}\ifdim\linewidth<5in\par\smallskip\else\hfill\fi
\begin{minipage}{\auditplotwidth}
\centering
\begin{tikzpicture}
\begin{axis}[
  width=0.95\linewidth,
  height=\auditplotheight,
  ybar=2pt,
  bar width=10pt,
  ymin=0,
  ymax=58,
  ytick={0,20,40},
  symbolic x coords={Pizza,Tacos,Salad,Pasta},
  xtick=data,
  enlarge x limits=0.18,
  title={\textbf{Graph B}},
  xlabel={Favorite lunch option},
  ylabel={Number of students},
  x tick label style={rotate=30,anchor=east,font=\scriptsize},
  tick label style={font=\scriptsize},
  label style={font=\scriptsize},
  title style={font=\small},
  ymajorgrids,
  nodes near coords,
  every node near coord/.append style={font=\scriptsize},
  axis on top
]
\addplot[fill=calloutblue,draw=dayblue] coordinates
  {(Pizza,52) (Tacos,47) (Salad,44) (Pasta,41)};
\end{axis}
\end{tikzpicture}
\end{minipage}
\end{center}
\begin{firsttakebox}{FIRST TAKE (5 min)}
Is pizza clearly the favorite? Name the graph that most influenced you.\par
\answer{Not graded. Expect Graph A to attract votes because its Pizza bar dominates.}
\end{firsttakebox}

\textit{Rules callout ``GRAPHING ONE CATEGORICAL VARIABLE (keep on the page)'' is printed on the student sheet.}\par\medskip

\dokbadge{1}~\textbf{(a)} Read the Pizza and Pasta counts. How many more students chose Pizza?\par
\answer{Pizza has 52 students and Pasta 41, so Pizza has 11 more.}

\dokbadge{2}~\textbf{(b)} Graph A starts at 40. Find the visible heights above 40 for Pizza and Pasta. Compare those heights with the actual counts.\par
\answer{Visible heights: Pizza 52-40=12 and Pasta 41-40=1. Pizza looks 12 times as tall, while the actual counts are 52 and 41, a difference of 11 students.}

\dokbadge{3}~\textbf{(c)} Which graph gives a fairer picture of the lunch votes? Explain how its scale affects the comparison. Use your counts to judge both `landslide' and `basically a tie,' then write a more accurate headline.\par
\answer{Graph B uses zero, so full bar heights represent counts. Graph A starts at 40 and exaggerates the difference: visible heights are 12 to 1, but counts are 52 to 41. Pizza has the most votes, yet neither a huge landslide nor an exact tie describes the counts. A supported headline is `Pizza leads Pasta by 11 votes; all four options receive support.' Accept other measured wording supported by the counts.}

\begin{summaryexitbox}{EXIT}
One thing the video changed about my first take: \hrulefill
\end{summaryexitbox}

\end{document}
